%! suppress = Makeatletter
%! suppress = TooLargeSection
%! suppress = MissingLabel
\documentclass{article}

% Fields
\usepackage{geometry}
\geometry{top=25mm}
\geometry{bottom=35mm}
\geometry{left=20mm}
\geometry{right=20mm}
% ------------------------------------------------

% Graphics
\usepackage{color}
\usepackage{tabularx}
\usepackage{tikz}
\usepackage{blkarray}
\usepackage{graphicx}
% ------------------------------------------------

% Math
\usepackage{amsmath, amsfonts}
\usepackage{amssymb}
\usepackage{proof}
\usepackage{mathrsfs}
% Crossed-out symbols
% https://tex.stackexchange.com/questions/75525/how-to-write-crossed-out-math-in-latex
\usepackage[makeroom]{cancel}
\usepackage{mathtools}
% ------------------------------------------------

% Additional font sizes
% https://www.overleaf.com/learn/latex/Questions/How_do_I_adjust_the_font_size%3F
\usepackage{moresize}
% Additional colors
% https://www.overleaf.com/learn/latex/Using_colours_in_LaTeX
\usepackage{xcolor}
% \texttimes
\usepackage{textcomp}
% ------------------------------------------------

% Language
\usepackage[utf8] {inputenc}
\usepackage[T2A] {fontenc}
\usepackage[english, russian] {babel}
\usepackage{indentfirst, verbatim}
\usetikzlibrary{cd, babel}
% ------------------------------------------------

% Fonts
\usepackage{stmaryrd}
\usepackage{cmbright}
\usepackage{wasysym}
% ------------------------------------------------

% Code
% https://tex.stackexchange.com/questions/99475/how-to-invoke-latex-with-the-shell-escape-flag-in-texstudio-former-texmakerx
% Colored, requires --shell-escape compiling option
\usepackage{minted}
\setminted{xleftmargin=\parindent, autogobble, escapeinside=\#\#}
\usepackage{listings}
% ------------------------------------------------

% Custom envs
% https://tex.stackexchange.com/questions/371286/draw-a-horizontal-line-in-latex
\newenvironment{proof}{\subparagraph{\hspace{-1em}Решение:\newline}\-}{\par\noindent\rule{\textwidth}{0.4pt}}
% ------------------------------------------------

% Custom commands
\newcommand{\comb}[1]{\mathbf{#1}}
\newcommand{\step}{\rightsquigarrow}
\newcommand{\term}[1]{\mathbf{#1}}
\newcommand{\ap}{~}
\newcommand{\termdef}{\coloneqq}
\newcommand{\subst}[3]{\left[#2 \mapsto #3 \right] #1}
\newcommand{\eqbeta}{=_\beta}
\newcommand{\eqbetaeta}{=_{\beta\eta}}
\newcommand{\betastep}{\rightarrow_\beta}
\newcommand{\betasteps}{\twoheadrightarrow_\beta}
\newcommand{\eqeta}{=_\eta}
\newcommand{\tlist}[1]{\term{[}#1\term{]}} % list-term
\renewcommand{\emph}[1]{{\color{blue} \textit{#1}}}
% ------------------------------------------------

% Head
\usepackage{fancyhdr}
\usepackage{hyperref}
\pagestyle{fancy}
\fancyhead[R]{Ярошевский Илья}
\fancyhead[L]{ИТМО MSE, ФП 2024, Дз 1}
% ------------------------------------------------

% Numbering
% https://tex.stackexchange.com/questions/80113/hide-section-numbers-but-keep-numbering
\makeatletter
\renewcommand\thesubsection{Блок \@arabic\c@subsection.\hspace{-0.8em}}
\renewcommand\thesubsubsection{Задание \@arabic\c@subsection.\@arabic\c@subsubsection\hspace{-0.8em}}
% https://tex.stackexchange.com/questions/327689/numbering-subsubsections-with-letters
\renewcommand\theparagraph{\alph{paragraph})\hspace{-0.8em}}
% https://tex.stackexchange.com/questions/129208/numbering-paragraphs-in-latex
\setcounter{secnumdepth}{4}
\makeatother
% ------------------------------------------------

\begin{document}
    \section*{Дз 2. Больше лямбда-исчисления}

    Глобальный минимум: 9б.

    \subsection{Комбинатор рекурсии}

    Локальный минимум: 2б.

    \subsubsection{}

    Решите следующие уравнения на лямбда-термы и произведите проверку решений.

    \paragraph{(1б)}

    $\forall M\ldotp F~M \eqbeta M~F$

    \begin{proof}
        \begin{gather}
            F~M \eqbeta M~F \\
            F~M \eqbeta \underbrace{(\lambda rec~m\ldotp m~rec)}_{F'}~F~M \\
            F \eqbetaeta F'~F \\
            F \coloneqq Y~F' \eqbeta Y~\lambda rec~m\ldotp m~rec
        \end{gather}
        Проверка:
        \begin{align}
            F M & \eqbeta Y~F'~M \eqbeta (\lambda rec~m\ldotp m~rec)~(Y~F')~M \\
                & \eqbeta M~(Y~F') \eqbeta M~F
        \end{align}
    \end{proof}

    \paragraph{(1б)}

    $\forall M~N\ldotp F~M~N \eqbeta F~N~(M~F)$

    \begin{proof}
        \begin{gather}
            F~M~N \eqbeta F~N~(M~F) \\
            F~M~N \eqbeta \underbrace{(\lambda rec~m~n\ldotp rec~n~(m~rec))}_{F'}~F~M~N \\
            F \eqbetaeta F'~F \\
            F \coloneqq Y~F' \eqbeta Y~\lambda rec~m~n\ldotp rec~n~(m~rec)
        \end{gather}
        Проверка:
        \begin{align}
            F~M~N & \eqbeta Y~F'~M~N \eqbeta (\lambda rec~m~n\ldotp rec~n~(m~rec))~(Y~F')~M~N \eqbeta \\
                  & \eqbeta (Y~F')~N~(M~(Y~F')) \eqbeta F~N~(M~F)
        \end{align}
    \end{proof}

    \subsubsection{(1б)}

    Даны термы $A$ и $B$.
    Найдите $F$ и $G$ такие, что
    \[
        \left\{
        \begin{aligned}
            F &\eqbeta A~F~G \\
            G &\eqbeta B~F~G \\
        \end{aligned}\right.
    \]

    Подсказка: сначала подумайте, что из чего нужно выражать.

    \begin{proof}
        \begin{gather}
            F \eqbeta A~F~G \eqbeta A~F~(B~F~G) \eqbeta \underbrace{(\lambda f\ldotp A~f~(B~f~G))}_{F'}~F \\
            F \coloneqq Y~F' \eqbeta Y~\lambda f\ldotp A~f~(B~f~G) \\
            G \eqbeta B~F~G \eqbeta B~(Y~\lambda f\ldotp A~f~(B~f~G))~G \eqbeta \underbrace{(\lambda g\ldotp B~(Y~\lambda f\ldotp A~f~(B~f~g))~g)}_{G'}~G \\
            G \coloneqq Y~G' \eqbeta Y~\lambda g\ldotp B~(Y~\lambda f\ldotp A~f~(B~f~g))~g
        \end{gather}
    \end{proof}

    \subsubsection{(2б)}

    Как мы уже увидели, в чистом лямбда-исчислении можно производить рекурсивные вычисления.
    И это при том, что в языке есть только переменные, аппликация и абстракция!
    Этого достигают с помощью построения комбинатора рекурсии, $\comb{Y}$ или других.

    В вашем любимом языке программирования рекурсия встроена неявно.
    Достаточно просто в теле функции вызвать её же саму по имени.
    Но что если имени нет, например, мы хотим заполучить рекурсивную лямбда-функцию?

    Реализуйте комбинатор рекурсии в вашем любимом языке программирования, используя встроенную в язык рекурсию.
    Ваш язык энергичный, предпримите меры от расходимости.

    С помощью написанного комбинатора реализуйте простейший рекурсивный факториал в виде \\ \texttt{factorial = \ldots}.

    За написание в статически-типизированном языке полагается дополнительный балл (но это может быть не просто).

    Подсказка: воспользуйтесь свойством комбинатора рекурсии: $\comb{Y}~F~M \eqbeta F~(\comb{Y}~F)~M$.

    \begin{proof}
    См. след. страницу
    \pagebreak
\begin{minted}[frame=lines,linenos=true,mathescape]{rust}
fn fix<'a, 'b: 'a, T, R>(
    f: &'b impl Fn(Box<dyn Fn(T) -> R + 'a>, T) -> R,
) -> Box<dyn Fn(T) -> R + 'a> {
    return Box::new(move |x| f(fix(f), x));
}

fn main() {
    let factorial: Box<dyn Fn(usize) -> u64> = fix(&|f, x| {
        if x == 0 {
            1
        } else {
            f(x - 1) * (x as u64)
        }
    });

    assert!(factorial(0) == 1);
    assert!(factorial(1) == 1);
    assert!(factorial(2) == 2);
    assert!(factorial(3) == 6);
    assert!(factorial(10) == 3628800);

    println!("Success");
}
\end{minted}
    \end{proof}

    \subsection{Программирование в $\lambda$-исчислении: числа и списки по Чёрчу}

    Локальный минимум: 4б.

    \subsubsection{}

    Задайте \emph{сильно нормализуемые}\footnote{На практике это означает, что нельзя использовать, например, $\comb{Y}$-комбинатор, но можно давать ответ в виде термов не в $\beta$-нормальной форме.} лямбда-термы с описанными свойствами.

    Подсказка: нетривиальный тут только первый пункт.

    \paragraph{(1б)}

    $\term{minus}$: если $n$ и $m$ ~--- числа Чёрча, то $\term{minus}~n~m = k$,
    если $\term{plus}~m~k = n$, и $\term{minus}~n~m = \term{0}$, если такого $k$
    нет.

    \begin{proof}
        \(\term{pred}\) из предыдущего ДЗ.
        \begin{gather}
            \term{pred} \coloneqq \lambda n\ldotp \term{fst}~(n~(\lambda p\ldotp \term{pair}~(\term{snd}~p)~(\term{suc}~(\term{snd}~p)))~(\term{pair}~\term{0}~\term{0})) \\
            \term{minus}~n~m \coloneqq m~(\lambda i\ldotp \term{pred}~i)~n
        \end{gather}
        Проверка:
        \begin{align}
            \term{minus}~2~1 & \eqbeta 1~(\lambda i\ldotp \term{pred}~i)~2 \eqbeta (\lambda i\ldotp \term{pred}~i)~2 \eqbeta \\
                             & \eqbeta \term{pred}~2 \eqbeta 1
        \end{align}
        \begin{align}
            \term{minus}~2~1 & \eqbeta 2~(\lambda i\ldotp \term{pred}~i)~1 \eqbeta (\lambda i\ldotp \term{pred}~i)~((\lambda i\ldotp \term{pred}~i)~1) \eqbeta \\
                             & \eqbeta (\lambda i\ldotp \term{pred}~i)~(\term{pred} 1) \eqbeta (\lambda i\ldotp \term{pred}~i)~0 \eqbeta \\
                             & \eqbeta \term{pred}~0 \eqbeta 0
        \end{align}
    \end{proof}

    \paragraph{(0.5б)}

    $\term{lt}$: если $n$ и $m$ ~--- числа Чёрча, то $\term{lt}~n~m = \term{true}$, если $n < m$, и $\term{false}$ в противном случае.

    \begin{proof}
        \begin{gather}
            \term{lt}~n~m \coloneqq \term{not}~(\term{isZero}~(\term{minus}~m~n))
        \end{gather}
        Проверка:
        \begin{align}
            \term{lt}~2~1 \eqbeta \term{not}~(\term{isZero}~(\term{minus}~1~2)) \eqbeta \term{not}~(\term{isZero}~0) \eqbeta \term{false}
        \end{align}
        \begin{align}
            \term{lt}~1~1 \eqbeta \term{not}~(\term{isZero}~(\term{minus}~1~1)) \eqbeta \term{not}~(\term{isZero}~0) \eqbeta \term{false}
        \end{align}
        \begin{align}
            \term{lt}~1~2 \eqbeta \term{not}~(\term{isZero}~(\term{minus}~2~1)) \eqbeta \term{not}~(\term{isZero}~1) \eqbeta \term{true}
        \end{align}
    \end{proof}

    \paragraph{(0.5б)}

    $\term{equal}$: если $n$ и $m$ ~--- числа Чёрча, то $\term{equal}~n~m = \term{true}$, если $n = m$, и $\term{false}$ в противном случае.

    \begin{proof}
        \begin{gather}
            \term{equal}~n~m \coloneqq \term{and}~(\term{isZero}~(\term{minus}~m~n))~(\term{isZero}~(\term{minus}~n~m))
        \end{gather}
        Проверка:
        \begin{align}
            \term{equal}~2~1 & \eqbeta \term{and}~(\term{isZero}~(\term{minus}~1~2))~(\term{isZero}~(\term{minus}~2~1)) \eqbeta \\
                             & \eqbeta \term{and}~(\term{isZero}~0)~(\term{isZero}~1) \eqbeta \term{false}
        \end{align}
        \begin{align}
            \term{equal}~2~2 & \eqbeta \term{and}~(\term{isZero}~(\term{minus}~2~2))~(\term{isZero}~(\term{minus}~2~2)) \eqbeta \\
                             & \eqbeta \term{and}~(\term{isZero}~0)~(\term{isZero}~0) \eqbeta \term{true}
        \end{align}
    \end{proof}

    \subsubsection{(1б)}

    Используя комбинатор примитивной рекурсии, напишите терм, который
    находит $n$-ое число Фибоначчи.

    \begin{proof}
        \begin{gather}
            \term{fibIter}~i~(\term{pair}~x~y) \coloneqq \term{pair}~(\term{plus}~x~y)~x  \\
            \term{fib}~n \coloneqq \term{snd}~(\term{rec}~\term{fibIter}~(\term{pair}~0~1)~n)
        \end{gather}
        \begin{align}
            \term{fib}~1 & \eqbeta \term{snd}~(\term{fibIter}~0~(\term{pair}~0~1)) \eqbeta \\
                         & \eqbeta \term{snd}~(\term{pair}~(\term{plus}~0~1)~1) \eqbeta \\
                         & \eqbeta \term{snd}~(\term{pair}~1~1) \eqbeta 1
        \end{align}
        \begin{align}
            \term{fib}~2 & \eqbeta \term{snd}~(\term{fibIter}~1~(\term{fibIter}~0~(\term{pair}~0~1))) \eqbeta \\
                         & \eqbeta \term{snd}~(\term{fibIter}~1~(\term{pair}~(\term{plus}~0~1)~1)) \eqbeta \\
                         & \eqbeta \term{snd}~(\term{fibIter}~1~(\term{pair}~1~1)) \eqbeta \\
                         & \eqbeta \term{snd}~(\term{pair}~(\term{plus}~1~1)~1) \eqbeta \\
                         & \eqbeta \term{snd}~(\term{pair}~2~1) \eqbeta 1
        \end{align}
        \begin{align}
            \term{fib}~4 & \eqbeta \term{snd}~(\term{fibIter}~3~(\term{fibIter}~2~(\term{fibIter}~1~(\term{fibIter}~0~(\term{pair}~0~1))))) \eqbeta \\
                         & \eqbeta \term{snd}~(\term{fibIter}~3~(\term{fibIter}~2~(\term{pair}~2~1))) \eqbeta \\
                         & \eqbeta \term{snd}~(\term{fibIter}~3~(\term{pair}~(\term{plus}~2~1)~2)) \eqbeta \\
                         & \eqbeta \term{snd}~(\term{fibIter}~3~(\term{pair}~3~2)) \eqbeta \\
                         & \eqbeta \term{snd}~(\term{pair}~(\term{plus}~3~2)~3) \eqbeta \\
                         & \eqbeta \term{snd}~(\term{pair}~5~3) \eqbeta 3
        \end{align}
    \end{proof}

    \subsubsection{}

    Задайте следующие простые сильно нормализуемые функции над списками.

    Здесь и далее подразумевается, что списки представлены в виде \textbf{списков Чёрча}.

    \paragraph{(0.5б)}

    Напишите функцию $\term{sum}$, суммирующий элементы списка: $\term{sum} \ap \tlist{\term{1}, \term{2}, \term{3}} \eqbeta \term{6}$.

    \begin{proof}
        \begin{gather}
            \term{sum}~xs \coloneqq \term{fold}~xs~\term{plus}~0
        \end{gather}
        Проверка:
        \begin{align}
            \term{sum} [1, 2, 3] \eqbeta \term{plus}~1~(\term{plus}~2~(\term{plus}~3~0)) \eqbeta
        \end{align}
    \end{proof}

    \paragraph{(0.5б)}

    Напишите функцию $\term{length}$, вычисляющий длину списка: $\term{length} \ap \tlist{\term{1}, \term{2}} \eqbeta \term{2}$

    \begin{proof}
        \begin{gather}
            \term{length}~xs \coloneqq \term{fold}~xs~(\lambda x~l\ldotp \term{succ}~l)~0
        \end{gather}
        Проверка:
        \begin{align}
            \term{length}~[3, 2, 1] & \eqbeta (\lambda i~l\ldotp \term{succ}~l)~3~(\term{fold}~[2, 1]~(\lambda i~l\ldotp \term{succ}~l)~0)) \eqbeta \\
                                    & \eqbeta \term{succ}~(\term{fold}~[2, 1]~(\lambda i~l\ldotp \term{succ}~l)~0)) \eqbeta \\
                                    & \eqbeta \term{succ}~(\term{succ}~(\term{succ}~0)) \eqbeta 3
        \end{align}
    \end{proof}

    \subsubsection{(1б)}

    Напишите функцию, которая сильно нормализуемо возвращает последний элемент непустого списка.

    \begin{proof}
        \begin{gather}
            \term{last}~xs \coloneqq \term{snd}~(\term{fold}~xs~(\lambda x~p\ldotp \term{if}~(\term{fst}~p)~p~(\term{pair}~\term{true}~x))~(\term{pair}~\term{false}~0))
        \end{gather}
        Проверка:
        \begin{align}
            \term{last}~[1, 2, 3] & \eqbeta \term{snd}~(\term{fold}~xs~(\lambda x~p\ldotp \term{if}~(\term{fst}~p)~p~(\term{pair}~\term{true}~x))~(\term{pair}~\term{false}~0)) \eqbeta \\
                                  % & \eqbeta \term{snd}~((\lambda p\ldotp \term{if}~(\term{fst}~p)~p~(\term{pair}~\term{true}~1))~((\lambda p\ldotp \term{if}~(\term{fst}~p)~p~(\term{pair}~\term{true}~2))~((\lambda p\ldotp \term{if}~(\term{fst}~p)~p~(\term{pair}~\term{true}~3))~(\term{pair}~\term{false}~0))))
                                  & \eqbeta \term{snd}~((\dots)~((\dots)~((\lambda p\ldotp \term{if}~(\term{fst}~p)~p~(\term{pair}~\term{true}~3))~(\term{pair}~\term{false}~0)))) \eqbeta \\
                                  & \eqbeta \term{snd}~((\dots)~((\dots)~(\lambda p\ldotp \term{if}~\term{false}~(\term{pair}~\term{false}~0)~(\term{pair}~\term{true}~3)))) \eqbeta \\
                                  & \eqbeta \term{snd}~((\dots)~((\lambda p\ldotp \term{if}~(\term{fst}~p)~p~(\term{pair}~\term{true}~2))~(\term{pair}~\term{true}~3))) \eqbeta \\
                                  & \eqbeta \term{snd}~((\dots)~(\lambda p\ldotp \term{if}~\term{true}~(\term{pair}~\term{true}~3)~(\term{pair}~\term{true}~2))) \eqbeta \\
                                  & \eqbeta \term{snd}~((\lambda p\ldotp \term{if}~(\term{fst}~p)~p~(\term{pair}~\term{true}~1))~(\term{pair}~\term{true}~3)) \eqbeta \\
                                  & \eqbeta \term{snd}~(\term{pair}~\term{true}~3) \eqbeta 3
        \end{align}
    \end{proof}

    \subsubsection{(2б)}

    Напишите функцию $\term{tail}$, которая сильно нормализуемо возвращает хвост списка, то есть список без первого элемента: например, $\term{tail}~\tlist{~} = \tlist{~}$ и $\term{tail}~\tlist{\term{5},\term{10},\term{13}} = \tlist{\term{10}, \term{13}}$.

    \begin{proof}
        \begin{gather}
            \term{tail}~xs \coloneqq \term{snd}~(\term{fold}~xs~(\lambda x~p\ldotp \term{pair}~(\term{cons}~x~(\term{fst}~p))~(\term{fst}~p))~(\term{pair}~[]~[]))
        \end{gather}
        Проверка:
        \begin{align}
            \term{tail}~[1, 2, 3] & \eqbeta \term{snd}~(\term{fold}~[1, 2, 3]~(\lambda x~p\ldotp \term{pair}~(\term{cons}~x~(\term{fst}~p))~(\term{fst}~p))~(\term{pair}~[]~[])) \eqbeta \\
            % \term{tail}~[1, 2, 3] \eqbeta \term{snd}~((\lambda p\ldotp \term{pair}~(\term{cons}~1~(\term{fst}~p))~(\term{fst}~p))~((\lambda p\ldotp \term{pair}~(\term{cons}~2~(\term{fst}~p))~(\term{fst}~p))~((\lambda p\ldotp \term{pair}~(\term{cons}~3~(\term{fst}~p))~(\term{fst}~p))~(\term{pair}~[]~[]))))
                                  & \eqbeta \term{snd}~((\dots)~((\dots)~((\lambda p\ldotp \term{pair}~(\term{cons}~3~(\term{fst}~p))~(\term{fst}~p))~(\term{pair}~[]~[])))) \eqbeta \\
                                  & \eqbeta \term{snd}~((\dots)~((\dots)~(\term{pair}~(\term{cons}~3~[])~[]))) \eqbeta \\
                                  & \eqbeta \term{snd}~((\dots)~((\lambda p\ldotp \term{pair}~(\term{cons}~2~(\term{fst}~p))~(\term{fst}~p))~(\term{pair}~[3]~[]))) \eqbeta \\
                                  & \eqbeta \term{snd}~(\term{pair}~[1, 2, 3]~[2, 3]) \eqbeta [2, 3] \\
        \end{align}
    \end{proof}

    \subsubsection{(1.5б)}

    Вспомним определение умножения чисел Чёрча:
    \[\term{mult} = \lambda n~m~s\ldotp n~(m~s)\]
    Рассмотрим его аналог для списков:
    \[\term{listMult} = \lambda n~m~s\ldotp n~(\lambda e\ldotp m~(\term{curry}~s~e))\]
    Что делает эта функция?
    Проредуцируйте наглядный пример для проверки предположения.

    \begin{proof}
        Это декартово произведение списков. То есть получаем список из всех возможных пар \(\langle a, b \rangle\), таких что \(a \in n, b \in m\).
        \begin{align}
            \term{listMult}~[1, 2]~[4, 5, 6] & \eqbeta \lambda s~z\ldotp ([4, 5, 6]~(\term{curry}~s~1))~(([4, 5, 6]~(\term{curry}~s~2))~z) \eqbeta \\
                                             & \eqbeta \lambda s~z\ldotp (\dots)~((\lambda z\ldotp [4, 5, 6]~((\lambda f~x~y\ldotp f~(\term{pair}~x~y))~s~2))~z) \eqbeta \\
                                             & \eqbeta \lambda s~z\ldotp (\dots)~(([4, 5, 6]~(\lambda y\ldotp s~(\term{pair}~2~y)))~z) \eqbeta \\
                                             & \eqbeta \lambda s~z\ldotp (\dots)~(s~(\term{pair}~2~4)~(s~(\term{pair}~2~5)~(s~(\term{pair}~2~6)~z))) \eqbeta \\
                                             & \eqbeta \lambda s~z\ldotp s~(\term{pair}~1~4)~(s~(\term{pair}~1~5)~(s~(\term{pair}~1~6)~(\dots)))) \eqbeta \\
                                             & \eqbeta [\term{pair}~1~4, \term{pair}~1~5, \term{pair}~1~6, \term{pair}~2~4, \term{pair}~2~5, \term{pair}~2~6]
        \end{align}
    \end{proof}

    \subsection{Стратегии редукции и нормальные формы}

    Локальный минимум: 1б.

    \subsubsection{(1б)}

    Дан терм:
    \[\lambda y\ldotp z~(\lambda x\ldotp \comb{I}~(\comb{I}~\comb{I})~(\comb{I}~\comb{I}))~x~y\]

    Дважды проредуцируйте его до нормальной формы, расписывая каждый шаг редукции: один раз с использованием аппликативной стратегии, во второй ~--- с
    использованием нормальной стратегии.

    \begin{proof}
        Аппликативная:
        \begin{align}
            & \lambda y\ldotp z~(\lambda x\ldotp \comb{I}~(\comb{I}~\comb{I})~(\comb{I}~\comb{I}))~x~y \betastep \\
            & \betastep \lambda y\ldotp z~(\lambda x\ldotp \comb{I}~\comb{I}~(\comb{I}~\comb{I}))~x~y \betastep \\
            & \betastep \lambda y\ldotp z~(\lambda x\ldotp \comb{I}~(\comb{I}~\comb{I}))~x~y \betastep \\
            & \betastep \lambda y\ldotp z~(\lambda x\ldotp \comb{I}~\comb{I})~x~y \betastep \\
            & \betastep \lambda y\ldotp z~(\lambda x\ldotp \comb{I})~x~y
        \end{align}
        \begin{align}
            & \lambda y\ldotp z~(\lambda x\ldotp \comb{I}~(\comb{I}~\comb{I})~(\comb{I}~\comb{I}))~x~y \betastep \\
            & \betastep \lambda y\ldotp z~(\lambda x\ldotp (\comb{I}~\comb{I})~(\comb{I}~\comb{I}))~x~y \betastep \\
            & \betastep \lambda y\ldotp z~(\lambda x\ldotp \comb{I}~(\comb{I}~\comb{I}))~x~y \betastep \\
            & \betastep \lambda y\ldotp z~(\lambda x\ldotp (\comb{I}~\comb{I}))~x~y \betastep \\
            & \betastep \lambda y\ldotp z~(\lambda x\ldotp \comb{I})~x~y \betastep \\
        \end{align}
        Нормальная:
    \end{proof}

    \subsubsection{}

    Приведите примеры замкнутых термов (и объясните выбор).

    \paragraph{(0.5б)}

    Находится в WHNF, но не в HNF.

    \begin{proof}
        \(\lambda x\ldotp \comb{I}~x\). Очевидно что находится в WHNF (на верху абстракция). Не находится в HNF, т.к. \(\comb{I}~x\) -- редекс.
    \end{proof}

    \paragraph{(0.5б)}

    Находится в HNF, но не в NF.

    \begin{proof}
        \(\lambda x\ldotp x~(\comb{I}~\comb{I})\). Находится в HNF, т.к. нельзя проредуцировать аппликацию \(x\), т.к. это переменная. Не находится в NF т.к. есть редекс \(\comb{I}~\comb{I}\).
    \end{proof}

    \subsubsection{(3б)}

    Две формулировки одного и того же задания в студию:
    \begin{itemize}
        \item Пусть дан терм $M$.
        Покажите, как подобрать такой терм $N$, что $N$ находится в нормальной форме и $N~\comb{I} =_\beta M$.
        \item Более привычная математикам формулировка: докажите, что для каждого $M$ существует $N$ в нормальной форме такой, что $N~\comb{I} \eqbeta M$.
    \end{itemize}
    \begin{proof}
        Пусть мета-переменная $\textit{Id}$ --- новая переменная, которая не встречается в терме $M$. Очевидно что терм $M$ содержит конечное количество редексов. Рассмотрим какой-нибудь поддтерм $P$ терма $M$, который является редексом $P \equiv A~B$. Подставим в исходный терм $M$ вместо этого поддетрма терм $(\textit{Id}~A)~B$, который не является редексом, т.к. \textit{Id} это переменная, получим новый терм $M_0$. Получается что терм $M_0$ содержит на 1 меньше редексов. Повторим операцию, но уже с $M_0$, пока не получим терм $M'$ в котором нет редексов.

        Покажем что $N \coloneqq \lambda \textit{Id}\ldotp M'$ удовлетворяет $N~\term{I} \eqbeta M$:
        \[ N~\term{I} \equiv (\lambda \textit{Id}\ldotp M')~\term{I} \eqbeta [\textit{Id} \mapsto \term{I}]M' \]
        Рассмотрим все подставленные поддетермы $[\textit{Id} \mapsto \term{I}]((\textit{Id}~A)~B) \equiv (\term{I}~A)~B \eqbeta A~B$. Получается что $N~\term{I} \eqbeta [\textit{Id} \mapsto \term{I}]M' \eqbeta M$.
    \end{proof}

\end{document}
